\maketitle

\begin{abstract}
Serving a long agent conversation means storing a key--value (KV) cache whose
size grows with the transcript, and shrinking that cache means deciding which
positions to discard.  Existing selectors read the model's own behaviour:
accumulated attention, or attention within an observation window.  A chat
transcript, however, already carries a second signal that costs nothing to
compute---its message roles, tool contracts, and constraint markers---and it is
natural to ask whether that free signal is sufficient.  This paper answers
\emph{when}.

We measure the operating range of a structural retention prior as a
dose--response curve rather than inferring it from endpoints.  Sweeping the
protected-role mass of a controlled agent benchmark across eight levels
(5.2\%--92.2\%), crossed with three keep budgets, over 18{,}240 scored
generations on Qwen3-8B, Llama-3.1-8B and Qwen3-32B, we find a sharp boundary: structural
retention is at ceiling while the protected set fits the budget and degrades
immediately once the oversubscription ratio $|S\cup M|/B$ exceeds
approximately~1.1.  The boundary replicates across all three model families, and a
fourfold increase in parameters reproduces the 8B curve to three significant
figures at every level, which is the signature of a property of the retention
policy rather than of the model.  Past
it, the ordering against attention-based selection reverses: structure loses to
matched-budget SnapKV in 20 of 24 Llama cells, winning no discordant pair in any
of them.

Pooling across task types conceals the mechanism, and the exception identifies
it.  At 18$\times$ oversubscription, tool-schema conformance remains at 1.000
while instruction supersession falls to exactly 0.000.  Because the evaluated
policy resolves ties by position, state carried in the leading system block
survives arbitrary oversubscription, whereas state accumulated across
mid-conversation turns does not.  The governing quantity is therefore not
protected mass alone but its interaction with where load-bearing state sits
relative to the tiebreak---a distinction that predicts, rather than
retrospectively excuses, the failure of the same prior on externally authored
agent traces.

We use that mechanism to make a falsifiable external prediction.  A blended
policy whose single weight is read from the prompt with nothing fitted never
loses to SnapKV across 48 cells and significantly exceeds it on Qwen at a 2\%
budget ($p = 4.9\times10^{-4}$), but the same mechanism predicts no benefit on
the Berkeley Function Calling Leaderboard, whose failures are dominated by loss
of multi-turn state rather than of schemas.  We pre-registered that prediction,
and the test confirms it: over 233 paired conversations ADAPT is
indistinguishable from SnapKV ($\Delta=-0.026$, $p=0.33$), inside the bound the
registration specified.  The method therefore confers no external benefit, while
the mechanism that predicts its failure does have external predictive power.  We
additionally report a matched-precision storage study as an explicit negative,
and document two experimental iterations that returned nulls, each for an
identified cause.
\end{abstract}

\section{Introduction}
\label{sec:introduction}

Autoregressive inference stores the keys and values of every prior position so
that each decoding step can attend over the prefix without recomputation.  The
resulting cache grows linearly with sequence length and, in agent workloads
where transcripts accumulate over many turns, becomes the binding constraint on
serving capacity~\cite{vaswani2017attention,kwon2023pagedattention}.  Systems
respond by paging the cache, by restricting attention to a sliding window, by
evicting positions judged unimportant, or by storing them at reduced precision
\cite{xiao2024streamingllm,zhang2023h2o,li2024snapkv,cai2024pyramidkv,liu2024kivi}.

Eviction requires a criterion, and the criteria in use are derived from the
model.  H2O accumulates attention mass per position; SnapKV scores prefix
positions by attention received from an observation window near the end of the
prompt; PyramidKV varies the per-layer budget according to how information
funnels through depth.  These signals are informative precisely because they
reflect what the model is doing, but obtaining them is not free: they require
attention statistics that exist only after the relevant computation has been
performed.

\begin{figure}[!htbp]
  \centering
  \includegraphics[width=\linewidth]{agent_trace_failure_mode.png}
  \caption{\textbf{Equal storage cost, unequal failure cost.}  Two policies keep
  the same illustrative 14-of-26-token budget.  Hatched cells are evicted;
  callouts show the resulting failure modes.  The figure is schematic and
  reports no empirical counts.}
  \label{fig:failure-mode}
\end{figure}

A chat transcript carries a different signal, and this is the observation that
motivates the present work.  Its serialization already distinguishes a system
preamble from a user turn, a tool schema from a tool result, an instruction from
filler.  That structure is available from the prompt alone, before any
computation, at negligible cost.  If a retention rule built on it were adequate,
a serving system could make eviction decisions without paying for attention
statistics at all.  The question is not whether such a rule is better than
attention-based selection---we will show it generally is not---but whether there
exists a characterisable regime in which the cheap signal suffices, and whether
that regime can be recognised in advance.

Structural retention itself is established.  RoleKV manages agent KV blocks by
semantic role rather than recency, though losslessly, deciding placement in the
memory hierarchy rather than discarding~\cite{rolekv}; work on structural-role
bias allocates budget across document-structural categories and reports that the
achievable gain is bounded by role density, becoming indistinguishable from
SnapKV where the informative role is sparse~\cite{structrolebias}.  What is not
established is the behaviour at the opposite extreme.  Agent traces do not
suffer from too few protected tokens; they suffer from too many, because tool
schemas can occupy nearly the whole context.  Prior work in this specific line,
including an earlier version of this study, asserted a boundary there on the
basis of two observations: a synthetic benchmark on which structural retention
performed well and carried little protected content, and an external benchmark
on which it failed and carried a great deal.  Two points do not establish a
curve, and a boundary inferred from the endpoints of an interval cannot
distinguish a real threshold from a coincidence of two unrelated effects.

\paragraph{Thesis.}  This paper replaces that inference with a measurement, and
in doing so arrives at a mechanism different from the one previously proposed.
Structural retention has a sharp operating boundary located where the protected
set ceases to fit the keep budget; the boundary is reproducible across model
families and computable before generation.  But the ratio alone does not
determine what is lost when the boundary is crossed.  Because the evaluated
policy breaks ties by position, state carried in the leading system block is
retained at any degree of oversubscription while state accumulated across
mid-conversation turns is discarded entirely.  It is this second factor, not the
ratio, that determines whether a workload is served well by the cheap signal.

\paragraph{Contributions.}
\begin{enumerate}
  \item \textbf{A measured operating range.}  We introduce a sweep methodology
  that parameterises the protected-role mass of a controlled agent benchmark
  while holding the task, its scoring, and its context length fixed, and we
  report the resulting dose--response curve over eight mass levels, three keep
  budgets, six retention arms and two model families
  (Sections~\ref{sec:sweep-design}--\ref{sec:sweep-results}).
  \item \textbf{The mechanism.}  We dissociate task families to show that a
  position-ordered tiebreak, not protected mass, decides which state survives
  oversubscription, and we show the dissociation replicates across
  architectures (Section~\ref{sec:sweep-mechanism}).
  \item \textbf{A falsifiable external prediction, registered and confirmed.}
  We derive from the mechanism a directional prediction for an externally
  authored benchmark and pre-register it---with a budget-selection rule and a
  power analysis---before the experiment is scored; the test then returns the
  predicted null ($\Delta=-0.026$, $p=0.33$, $n=233$ paired), establishing
  external predictive power for the mechanism while confirming that the method
  itself does not transfer (Section~\ref{sec:external}).
  \item \textbf{Negative and null results reported in full.}  We report a
  storage study that finds no quality separation, two sweep iterations that
  returned nulls with identified causes, and a failure-mode census showing that
  the external benchmark is limited more by model capability than by retention
  (Sections~\ref{sec:storage}, \ref{sec:sweep-design}, \ref{sec:external-census}).
\end{enumerate}

\section{Background and Related Work}
\label{sec:related}

\paragraph{Paging and streaming.}
PagedAttention separates logical KV blocks from non-contiguous physical storage
to reduce fragmentation and enable sharing across
requests~\cite{kwon2023pagedattention}; it manages placement rather than
selecting which positions deserve to remain.  StreamingLLM observes that a small
number of initial tokens act as attention sinks and that retaining them
alongside a recent window supports unbounded
streaming~\cite{xiao2024streamingllm}.  We adopt sink-plus-recent as a mandatory
region for every policy and as the shape of our position-only control, without
claiming to reproduce the streaming setting in which it was introduced.

\paragraph{Attention-derived eviction.}
H2O retains heavy hitters by accumulated attention together with recent
positions~\cite{zhang2023h2o}.  SnapKV selects prefix positions using attention
from an observation window at the end of the prompt, exploiting the empirical
stability of those patterns~\cite{li2024snapkv}.  PyramidKV allocates budget
non-uniformly across layers~\cite{cai2024pyramidkv}.  All three infer importance
from model behaviour.  We use these methods as the standard against which the
sufficiency of a behaviour-free signal is judged.  Except where noted, the
attention comparison in this paper uses the \texttt{kvpress} implementation of
SnapKV rather than a reimplementation.

\paragraph{Role- and structure-aware retention.}
The idea that a token's structural role should inform its retention is not new,
and two lines of prior work bear directly on ours.  RoleKV tags KV blocks with
agent-workflow roles---system, instruction, reasoning, action, tool
output---and manages them by role rather than recency, motivated by an inverted
age--importance effect in which the oldest blocks remain the most
attended~\cite{rolekv}.  RoleKV is, however, a \emph{lossless} policy: it
decides where blocks reside in the memory hierarchy and does not alter model
outputs.  The present work studies lossy eviction, where a discarded position
cannot be recovered and the failure modes are correspondingly different.

Closer to our question, work on structural-role bias applies a role-conditional
budget allocator over document-structural categories such as schema labels,
table headers and prose, and reports that the achievable gain is \emph{bounded
by role density}: where the informative role occupies a small fraction of the
context, the policy becomes indistinguishable from
SnapKV~\cite{structrolebias}.  That result establishes a boundary at the
\emph{sparse} end of the role-density axis.  Our contribution is the opposite
end.  We show that when protected content instead saturates the budget, the
policy does not merely lose its advantage but inverts its ordering against
attention-based selection, and---the finding we regard as
new---that \emph{which} state is lost is determined not by density but by the
position-ordered tiebreak the policy falls back on
(Section~\ref{sec:sweep-mechanism}).  The two results are complementary
descriptions of when a structural prior is usable.

\paragraph{Diagnostics for when eviction helps.}
A related line asks not which selector is best but when eviction succeeds at
all.  A fixed-contract diagnostic decomposes eviction failure into access
estimation, value ranking, and score-to-cache projection, and derives a
pre-generation margin that predicts whether value-aware selection will
help~\cite{valuediag}.  That study evaluates the same two 8B models used here
and explicitly confines itself to attention- and value-aware selection, leaving
role-based and structural retention unexamined.  The diagnostic we report---the
oversubscription ratio together with the location of load-bearing state---is
the structural counterpart it does not cover.

\paragraph{Quantization.}
KIVI shows that asymmetric group-wise low-bit quantization of the KV cache
preserves quality while keeping a residual of recent tokens at full
precision~\cite{liu2024kivi}.  Our storage study is not a new quantizer and is
not positioned against KIVI; it asks the narrower question of whether
\emph{role-aware placement} of a fixed INT4 fraction outperforms uniform
placement, and answers in the negative.  FlashInfer supplies the composable
kernels and log-sum-exp state merging that our mixed-precision decode path
uses~\cite{ye2025flashinfer}.

\paragraph{Evaluation.}
LongBench and RULER assess long-context capability
broadly~\cite{bai2024longbench,hsieh2024ruler}; our benchmark instead isolates
failure attribution in synthetic agent protocols and cannot substitute for
them.  BFCL provides stateful multi-turn function-calling tasks with an
executable checker~\cite{patil2025bfcl}, and $\tau$-bench evaluates
policy-constrained tool interaction with a simulated user~\cite{yao2024taubench}.
We use BFCL for external task success and public $\tau$-bench trajectories for a
generation-free retention audit; neither informed the construction of our
benchmark.

\section{Retention Policies and the Oversubscription Boundary}
\label{sec:method}

\subsection{Matched-budget retention}
\label{sec:retention-method}

Let a tokenised trace be $x_{1:n}$ with structural roles $r_{1:n}$ assigned by a
tagger that reads message roles and transparent schema or constraint markers.
The tagger inspects neither benchmark answers nor future attention.  For keep
fraction $k$ the nominal budget is
\begin{equation}
  B(n,k)=\operatorname{round}(kn),
  \label{eq:budget}
\end{equation}
and the mandatory set is $M=M_{\mathrm{sink}}\cup M_{\mathrm{recent}}$, holding
the first 16 and last 128 positions.  Every compressed arm selects a retained
set $A\supseteq M$ with
\begin{equation}
  |A|=\max\bigl(B(n,k),\,|M|\bigr),
  \label{eq:matched}
\end{equation}
so that arms differ only in \emph{which} positions survive, never in how many.
Position-blind retention fills from the recent edge.  Structural retention first
admits system, tool, constraint and other state-bearing positions, then fills
any residual budget from the recent edge.

\subsection{The boundary, and why the tiebreak matters}
\label{sec:boundary-theory}

Write $S$ for the structurally protected positions.  If $|S\cup M|\le B$, every
protected position fits and the binary role label separates protected content
from filler.  If $|S\cup M|>B$, the protected set is \emph{oversubscribed}: a
binary label cannot order positions within $S$, and the policy must resolve the
excess by some other rule.  We define the oversubscription ratio
\begin{equation}
  \rho \;=\; \frac{|S\cup M|}{B},
  \label{eq:ratio}
\end{equation}
which is computable from the prompt before any token is generated.

The literature framing, and our own earlier one, stops here and treats $\rho>1$
as the condition under which the prior ceases to be useful.  That is incomplete.
When $\rho>1$ the policy does not fail arbitrarily; it falls back on whatever
secondary ordering it implements.  The evaluated policy admits protected
positions in index order, so the surviving subset of $S$ is its \emph{earliest}
elements.  Whether the trace's load-bearing state is inside that surviving
prefix is therefore a second, independent condition.  Section~\ref{sec:sweep}
shows that this second condition, not $\rho$, decides which tasks break.

\subsection{ADAPT: a budget-relative blend}
\label{sec:adapt-method}

A hard union of structural protection with attention-based selection consumes
the whole budget whenever $\rho>1$, and in earlier experiments performed no
better than SnapKV alone and worse than both parents at tight budgets.  ADAPT
replaces the hard constraint with a budget-relative prior,
\begin{equation}
  \alpha=\min\!\left(1,\frac{B}{|S|}\right),\qquad
  s=\alpha\cdot\mathrm{rank}(s_{\text{struct}})+(1-\alpha)\cdot
  \mathrm{rank}(s_{\text{attn}}),
  \label{eq:adapt}
\end{equation}
where both quantities are known from the prompt and no scalar is fitted to any
outcome.  Rank normalisation is required because the structural bands are of
order $10^{6}$ while attention scores are of order $10^{-3}$; a raw weighted sum
would be dominated by scale rather than by preference.  At $\alpha=1$ ADAPT
recovers the structural ordering exactly; as oversubscription grows it shifts
weight toward attention.  ADAPT computes attention scores and is therefore
strictly more expensive than SnapKV; it is evaluated for accuracy at aggressive
compression, not as a cost reduction.

\section{PriorityBench-A}
\label{sec:benchmark}

PriorityBench-A generates multi-turn chats in which a designated span carries
the state required by a final request, while intervening turns supply benign
filler.  Scoring is binary and programmatic, with no model judge.  Three task
families are represented equally (Table~\ref{tab:tasks}).  Two frozen manifests
serve different studies: a 240-example set spanning 8k, 16k and 32k nominal
contexts for the storage study, and a 120-example stress set at 8k and 16k for
retention.

\begin{table}[!htbp]
  \centering
  \caption{PriorityBench-A task families.  All scorers are deterministic.}
  \label{tab:tasks}
  \begin{tabular}{lll}
    \toprule
    Family & Required behaviour & Scorer \\
    \midrule
    Tool schema & Emit a valid call under an early contract & JSON/schema subset \\
    Instruction supersession & Follow the latest conflicting constraint & Required/forbidden regex \\
    Multi-turn state & Reuse an earlier identifier, path, or preference & Exact slot match \\
    \bottomrule
  \end{tabular}
\end{table}

Because the generator's templates are visible to the tagger, a heuristic could
appear semantic merely by recognising them.  We therefore evaluate two
transformations of each slice.  \emph{Middle relocation} moves state away from
the prefix while holding the leading system message and the final request fixed.
\emph{Buried state} lengthens or embeds state-bearing turns so that a
short-span heuristic cannot protect all of the gold.  Separately, a CPU audit
maps gold-bearing messages to tokenizer spans and intersects them with the
selected positions, using no model outputs, to measure how much gold already
lies within the always-kept window.  On the stress slices only 1.0--1.3\% of
Qwen gold tokens and none of the Llama gold tokens fall inside sink plus recent,
so the in-distribution gap over position-blind eviction reflects retention
rather than a labelling leak.

\subsection{In-distribution behaviour}
\label{sec:in-dist}

At a 25\% keep budget on the stress set, structural retention passes 112 of 120
Qwen examples against 1 of 120 for the position-blind control.  We report this
result but decline to build on it: the control retains only sink plus recent on
a benchmark whose gold state is deliberately placed outside that window, and is
therefore close to guaranteed to fail.  The informative comparison is against
attention-based selection, where structure passes 112/120 and SnapKV,
PyramidKV, and a structure--SnapKV hybrid each pass 108/120; the paired
structure--SnapKV difference is not significant (exact McNemar, $b=4$, $c=0$,
$p=0.125$).  Under middle relocation, structure and FullKV both reach 115/120;
under buried state, structure falls to 80/120 against FullKV at 107/120.  Visible
roles are thus a useful prior, not an oracle for arbitrarily embedded state.

\begin{figure}[!htbp]
  \centering
  \includegraphics[width=0.86\linewidth]{eviction_and_baselines.png}
  \caption{\textbf{In-distribution comparison at a 25\% budget}
  ($n=120$, Qwen3-8B).  Points are pass rates with Wilson 95\% intervals; the
  bracket reports the paired structure--SnapKV McNemar test.  H2O* denotes a
  chunked reimplementation carrying a fidelity caveat.  This operating point is
  the one at which structure and attention selection agree;
  Section~\ref{sec:sweep-results} shows the agreement does not survive
  oversubscription.}
  \label{fig:eviction}
\end{figure}

\section{Measuring the Operating Boundary}
\label{sec:sweep}

\subsection{Design}
\label{sec:sweep-design}

To convert the boundary from an inference into a measurement we require traces
whose protected-role mass varies continuously while everything else is held
constant.  We obtain them by converting filler user/assistant pairs into
\emph{distractor} tool-schema pairs: a short registry notice followed by a JSON
dump of deferred schemas, sized to the characters it replaces.  The gold task,
its scoring payload, and the approximate context length are unchanged; only the
composition of the surrounding context varies.  Eight levels span 5.2\%--92.2\%
measured protected-role mass, crossed with keep budgets
$k\in\{0.02,0.05,0.10\}$ at 16k nominal context, with $n=60$ examples per cell.
The same base example appears at every level, so comparisons along the axis are
paired on the task.  Every non-FullKV arm is a \texttt{kvpress} press over an
identical full prefill at an identical compression ratio, so the arms differ
only in policy and never in mechanism.

Two earlier iterations of this design returned nulls, and both are informative
enough to report.  In the first, gold state occupied its natural position in the
template prefix; every non-blind arm then scored 1.000 in all 24 cells,
including at ninefold oversubscription.  The cause is the tiebreak of
Section~\ref{sec:boundary-theory}: when the protected set does not fit, the
policy retains its earliest members, which in that layout are the gold positions
themselves.  Oversubscription could not bite regardless of how large it became.
In the second, gold was relocated to mid-context but the transform was applied
after schema conversion; because the relocation helper classifies converted
schema turns as non-filler, it bundled them with the gold block and inserted
that block after a shrinking filler prefix.  Gold consequently drifted back
toward the start as protected mass rose---from character fraction 0.65 at the
lowest level to 0.07 at the highest---cancelling the effect that the sweep was
designed to expose.  The reported sweep applies relocation \emph{before}
conversion, with gold verified at character fraction 0.63--0.65 at every level.

\subsection{A sharp boundary that replicates}
\label{sec:sweep-results}

Structural retention is at ceiling while the protected set fits the budget and
degrades immediately once it does not (Table~\ref{tab:sweep},
Figure~\ref{fig:sweep}).  The knee lies near $\rho\approx1.1$ on both model
families.  FullKV passes every cell, so the collapse reflects retention rather
than task difficulty, and the position-blind arms remain near zero throughout.

\begin{table}[!htbp]
  \centering
  \caption{Pass rate against oversubscription ratio $\rho$, $n=60$ per cell.
  Ranges span the cells falling in each band.}
  \label{tab:sweep}
  \begin{tabular}{lrrrrr}
    \toprule
    $\rho$ & Structure (Qwen) & Structure (Llama) & SnapKV & ADAPT & FullKV \\
    \midrule
    $0.5\times$ & 1.000 & 1.000 & 1.000 & 1.000 & 1.000 \\
    $1.0\times$ & 1.000 & 1.000 & 1.000 & 1.000 & 1.000 \\
    $1.1\times$ & 1.000 & 0.883 & 1.000 & 1.000 & 1.000 \\
    $2.2\times$ & 0.617 & 0.600 & 1.000 & 1.000 & 1.000 \\
    $4$--$18\times$ & 0.400 & 0.333--0.417 & 0.967--1.000 & 1.000 & 1.000 \\
    $27$--$45\times$ & 0.067 & 0.000 & 0.750--0.867 & 0.950--1.000 & 1.000 \\
    \bottomrule
  \end{tabular}
\end{table}

\begin{figure}[!htbp]
  \centering
  \includegraphics[width=0.49\linewidth]{sweep_boundary_qwen.png}\hfill
  \includegraphics[width=0.49\linewidth]{sweep_boundary_llama.png}
  \caption{\textbf{One boundary, two architectures.}  Pass rate against
  oversubscription ratio for Qwen3-8B (left) and Llama-3.1-8B (right).  Marker
  shape encodes the keep budget, so points originating from 2\%, 5\% and 10\%
  budgets can be seen to lie on a common curve.  The dotted line marks the
  protected set exactly filling the budget.}
  \label{fig:sweep}
\end{figure}

Past the boundary, the parity with attention-based selection reported in
Section~\ref{sec:in-dist} reverses.  On Llama, structure loses to matched-budget
SnapKV in 20 of 24 cells, with $b=0$ in every one---that is, structure wins no
discordant pair anywhere.  On Qwen the same pattern holds at the tightest
budget, with $b=0$ and $c$ between 41 and 54.  The in-distribution parity is
therefore a property of one operating point rather than a general result, and
should not be reported without it.

\subsection{The boundary is a property of the policy, not the model}
\label{sec:sweep-scale}

If the boundary reflected how a particular model attends, its location or depth
should move with model capacity.  We tested this by repeating the sweep on
Qwen3-32B, four times the parameter count of the primary matrix, sharded across
two GPUs, at a fixed $k=0.05$ so that protected mass alone varies $\rho$ from
1.0 to 18.2 ($n=60$ per cell, 2{,}880 work units, no failures).

\begin{table}[!htbp]
  \centering
  \caption{Structural retention at $k=0.05$ on Qwen3-32B against the 8B result
  at the same cells.  Paired tests compare structure with matched-budget SnapKV
  on the 32B run.}
  \label{tab:scale}
  \begin{tabular}{lrrrrl}
    \toprule
    $\rho$ & Structure (32B) & Structure (8B) & SnapKV (32B) & ADAPT (32B) & Structure vs SnapKV \\
    \midrule
    $1.0\times$  & 1.000 & 1.000 & 1.000 & 1.000 & $b{=}0$, $c{=}0$, $p=1.00$ \\
    $2.2\times$  & 0.617 & 0.617 & 1.000 & 1.000 & $b{=}0$, $c{=}23$, $p=2.4\times10^{-7}$ \\
    $3.7\times$  & 0.417 & 0.417 & 1.000 & 1.000 & $b{=}0$, $c{=}35$, $p=5.8\times10^{-11}$ \\
    $5.3\times$  & 0.400 & 0.400 & 1.000 & 1.000 & $b{=}0$, $c{=}36$, $p=2.9\times10^{-11}$ \\
    $8.0\times$  & 0.400 & 0.400 & 0.983 & 1.000 & $b{=}1$, $c{=}36$, $p=5.5\times10^{-10}$ \\
    $11.0\times$ & 0.433 & 0.400 & 1.000 & 1.000 & $b{=}0$, $c{=}34$, $p=1.2\times10^{-10}$ \\
    $14.5\times$ & 0.400 & 0.400 & 0.967 & 1.000 & $b{=}2$, $c{=}36$, $p=5.4\times10^{-9}$ \\
    $18.2\times$ & 0.400 & 0.400 & 0.983 & 1.000 & $b{=}1$, $c{=}36$, $p=5.5\times10^{-10}$ \\
    \bottomrule
  \end{tabular}
\end{table}

The curves agree to three significant figures at every level.  This is a
stronger statement than consistency: the knee does not move, the descent does
not soften, and the floor does not shift, despite a fourfold increase in
capacity.  Section~\ref{sec:sweep-mechanism} explains why it cannot.  The floor
of 0.400 is the tool-schema family passing while the other two fail; which
family passes is fixed by where the policy's tiebreak retains positions, and
that is arithmetic over the role assignment rather than behaviour of the
network.  A property of the retention policy is invariant to the model it is
applied to, and this is the prediction that invariance makes.

\subsection{Which state survives, and why}
\label{sec:sweep-mechanism}

Pooling across task families conceals the mechanism, and the concealment is not
subtle: the pooled floor near 0.400 is the average of two families at 1.000 and
0.000 respectively.  Table~\ref{tab:dissoc} separates them.

\begin{table}[!htbp]
  \centering
  \caption{Structural retention by task family at high oversubscription
  ($\rho\ge5$), $n=20$ per cell per model.}
  \label{tab:dissoc}
  \begin{tabular}{llrrr}
    \toprule
    Model & Budget & Tool schema & Supersession & Multi-turn state \\
    \midrule
    Qwen  & 0.10 & 1.000 & 0.000 & 0.200 \\
    Qwen  & 0.05 & 1.000 & 0.000 & 0.200 \\
    Qwen  & 0.02 & 0.000 & 0.000 & 0.200 \\
    Llama & 0.10 & 1.000 & 0.000 & 0.000--0.150 \\
    Llama & 0.05 & 1.000 & 0.000 & 0.000--0.150 \\
    Llama & 0.02 & 0.000 & 0.000 & 0.000--0.050 \\
    \bottomrule
  \end{tabular}
\end{table}

Tool-schema conformance survives eighteenfold oversubscription intact, while
instruction supersession is extinguished.  The explanation is structural rather
than statistical.  A tool-schema item's load-bearing content is the leading
system block; the policy admits protected positions in index order; therefore
the contract survives however oversubscribed the policy becomes, until the
budget is too small to hold the block at all---which is precisely the
$k=0.02$ row, where tool-schema conformance also collapses.  Supersession and
multi-turn state live in mid-conversation turns, which the same ordering
discards first.

This supersedes the account of Section~\ref{sec:boundary-theory} in an important
way.  The ratio $\rho$ governs \emph{whether} a binary structural label can still
discriminate.  The position of load-bearing state relative to the tiebreak
governs \emph{what is lost when it cannot}.  Only the pair is predictive, and a
serving policy consulting $\rho$ alone would be systematically over-optimistic
for exactly those workloads---agent traces with accumulating state---that
motivate structural retention in the first place.

\subsection{ADAPT across the range}
\label{sec:sweep-adapt}

ADAPT holds 0.950--1.000 on Qwen and 0.983--1.000 on Llama across all 24 cells
per model, and loses no discordant pair to SnapKV in any of the 48 cells.  Its
measured $\alpha$ decays monotonically from 1.000 to 0.022 as oversubscription
grows, with no value fitted to any outcome.

On Qwen at $k=0.02$ it significantly exceeds SnapKV in four cells
(Table~\ref{tab:adapt}), always with $c=0$.  The advantage does not replicate on
Llama, where SnapKV never falls below 0.983 at any tested setting and therefore
leaves no headroom to recover.  We claim only that ADAPT matches attention-based
selection everywhere tested and exceeds it where attention-based selection
itself begins to degrade.

\begin{table}[!htbp]
  \centering
  \caption{ADAPT versus SnapKV on Qwen3-8B at $k=0.02$, exact two-sided
  McNemar, $n=60$ per cell.}
  \label{tab:adapt}
  \begin{tabular}{lrrrrr}
    \toprule
    Protected mass & ADAPT & SnapKV & $b$ & $c$ & $p$ \\
    \midrule
    0.399 & 1.000 & 0.900 & 6  & 0 & $3.1\times10^{-2}$ \\
    0.549 & 1.000 & 0.867 & 8  & 0 & $7.8\times10^{-3}$ \\
    0.723 & 1.000 & 0.867 & 8  & 0 & $7.8\times10^{-3}$ \\
    0.909 & 0.950 & 0.750 & 12 & 0 & $4.9\times10^{-4}$ \\
    \bottomrule
  \end{tabular}
\end{table}

\section{External Evaluation}
\label{sec:external}

\subsection{The prior fails, as the mechanism predicts}
\label{sec:external-transfer}

BFCL V3 multi-turn provides an externally authored context distribution scored
by an unmodified official checker that executes tool calls against stateful API
implementations~\cite{patil2025bfcl}.  At a 25\% keep budget over 140 paired
Qwen conversations, FullKV reaches 0.193 and SnapKV 0.136, a difference that is
not significant ($p=0.152$); the hard structural policy scores 0.000, and the
paired comparison against FullKV gives $p=1.5\times10^{-8}$.  The same ordering
holds on 143 Llama conversations.  BFCL prompts are dominated by tool schemas
and carry 98.8\% protected-role mass, so any tight budget is oversubscribed by
more than an order of magnitude.

\begin{figure}[!htbp]
  \centering
  \includegraphics[width=0.86\linewidth]{external_bfcl_arms.png}
  \caption{\textbf{Paired BFCL accuracy with Wilson intervals.}  The ordering
  holds on both model families: no significant FullKV--SnapKV difference is
  detected, while the non-attention policies sit at zero with a Wilson upper
  bound of 0.027.}
  \label{fig:external-arms}
\end{figure}

\subsection{Why the workload cannot reward the prior}
\label{sec:external-census}

A census of failure modes on the frozen points explains the result more
precisely than the ratio does, and does so without new generation
(Table~\ref{tab:census}).

\begin{table}[!htbp]
  \centering
  \caption{BFCL failure modes by arm at $k=0.25$, official checker error types.}
  \label{tab:census}
  \begin{tabular}{lrrrr}
    \toprule
    Arm & Pass & Empty turn response & State mismatch & Execution mismatch \\
    \midrule
    FullKV    & 20.5\% & 45.5\% & 22.7\% & 11.4\% \\
    SnapKV    & 13.9\% & 50.4\% & 21.2\% & 14.6\% \\
    ADAPT     & 12.9\% & 46.4\% & 27.1\% & 13.6\% \\
    Structure & 0.0\%  & 78.0\% & 10.0\% & 12.0\% \\
    \bottomrule
  \end{tabular}
\end{table}

Two facts follow.  First, 45.5\% of conversations fail with an empty or
unparseable turn response \emph{even with the full cache}, and FullKV passes only
20.5\%; on this workload an 8B model is limited more by its own ability to emit
well-formed multi-turn tool calls than by retention.  The entire interval within
which a retention policy can be judged is the roughly eight conversations
separating SnapKV from FullKV.  Second, the residual failures are dominated by
state mismatch---loss of accumulated multi-turn state---which is precisely the
class that Section~\ref{sec:sweep-mechanism} shows the structural prior does not
protect.  The schemas that it does protect sit in the system block and are
evidently not the binding constraint.  Consistent with this, the ADAPT arm shows
a higher state-mismatch rate than SnapKV (27.1\% versus 21.2\%) at marginally
lower accuracy.

\subsection{A pre-registered test and its budget}
\label{sec:external-prereg}

The mechanism therefore yields a directional prediction, which we registered
before scoring any conversation at a tight budget: on BFCL, ADAPT should match
SnapKV within noise or fall slightly below it, because at this workload's
protected mass the structural component biases retention toward schema tokens at
the expense of the state the checker scores.  The registration also fixes the
analysis (exact McNemar paired on task, $n=300$), the interpretation of every
possible outcome, and a minimum detectable effect; with $c=0$ the design detects
$b\ge6$, and a null bounds any true advantage below roughly 0.05 absolute.

Selecting the budget required care, because the regime in which ADAPT can help
is one where attention-based selection is degraded but not extinguished.  A
pre-specified rule chose the tightest rung at which SnapKV falls between 0.25 and
0.70 of FullKV on the same conversations.  A ladder of SnapKV-only pilots
(Table~\ref{tab:ladder}) shows that attention-based selection floors on BFCL
below a 10\% budget, so the tight budgets at which the synthetic advantage
appears do not exist on this workload; the rule selects $k=0.15$.

\begin{table}[!htbp]
  \centering
  \caption{Budget-selection pilot, SnapKV only, against the budget-independent
  FullKV reference on the same conversations.}
  \label{tab:ladder}
  \begin{tabular}{lrrrl}
    \toprule
    Budget & $n$ & SnapKV & FullKV & SnapKV/FullKV \\
    \midrule
    0.02 & 49 & 0.000 & 0.184 & 0.00 \\
    0.05 & 49 & 0.000 & 0.184 & 0.00 \\
    0.10 & 48 & 0.021 & 0.188 & 0.11 \\
    0.15 & 89 & 0.124 & 0.202 & 0.61 \\
    0.20 & 92 & 0.087 & 0.207 & 0.42 \\
    \bottomrule
  \end{tabular}
\end{table}

\subsection{Outcome: the predicted null}
\label{sec:external-outcome}

The main test ran 300 conversations per arm at $k=0.15$, yielding 233 fully
paired across FullKV, SnapKV and ADAPT at 94\% paired completeness, with zero
matched-budget violations and 33 exclusions, all of them context-limit
exclusions distributed across arms (Table~\ref{tab:external-main}).

\begin{table}[!htbp]
  \centering
  \caption{BFCL V3 multi-turn at the pre-registered budget $k=0.15$, official
  checker, $n=233$ fully paired conversations.  Intervals are Wilson 95\%;
  tests are exact two-sided McNemar.}
  \label{tab:external-main}
  \begin{tabular}{lrrl}
    \toprule
    Arm & Passes & Accuracy & Wilson 95\% \\
    \midrule
    FullKV & 52/233 & 0.223 & [0.174, 0.281] \\
    SnapKV & 32/233 & 0.137 & [0.099, 0.187] \\
    ADAPT  & 26/233 & 0.112 & [0.077, 0.158] \\
    \midrule
    \multicolumn{4}{l}{\emph{Paired comparisons}} \\
    ADAPT vs SnapKV & $b{=}10$, $c{=}16$ & $\Delta=-0.026$ & $p=0.327$ \\
    SnapKV vs FullKV & $b{=}12$, $c{=}32$ & $\Delta=-0.086$ & $p=0.0037$ \\
    ADAPT vs FullKV & $b{=}12$, $c{=}38$ & $\Delta=-0.112$ & $p=0.0003$ \\
    \bottomrule
  \end{tabular}
\end{table}

ADAPT is statistically indistinguishable from SnapKV and trends 2.6 points
below it.  This is the registered prediction, and the observed effect lies well
inside the bound the registration specified: a null on this design constrains
any true ADAPT advantage to below roughly 0.05 absolute, and the measured
$\Delta$ is $-0.026$ with a 10/16 discordant split.  The mechanism of
Section~\ref{sec:sweep-mechanism} therefore made a quantitative, directional,
pre-registered prediction about a benchmark the authors did not construct, and
the prediction was borne out.

We emphasise what this does and does not establish.  It does not establish that
the method is useful on BFCL; it establishes the opposite, and the paper claims
no external benefit.  What it supports is the \emph{mechanism}: a structural
prior that protects system-block content cannot help a workload whose residual
failures are losses of accumulated multi-turn state, and the size and direction
of that shortfall were derivable in advance from a synthetic sweep.  The two
secondary comparisons are consistent with the same account.  SnapKV trails
FullKV significantly ($p=0.0037$), confirming that a 15\% budget is genuinely
lossy on this workload rather than free, and ADAPT trails FullKV more strongly
($p=0.0003$), as expected when 98.8\% protected mass oversubscribes the budget
more than sixfold.

\subsection{A generation-free corroboration}
\label{sec:external-tau}

An audit of 4{,}856 public $\tau$-bench trajectories (828{,}000 extracted spans),
using no model outputs, corroborates the mechanism from the retention side:
structural retention preserves explicit policy lines at 0.820 against 0.001 for
a position-only budget, but loses on reused identifiers (0.055 against 0.315).
The prior wins precisely where protected spans are scarce and separable, and
loses on the recently referenced values that a recency heuristic captures.

\section{Storage: A Matched-Precision Negative}
\label{sec:storage}

We tested separately whether role-awareness helps when positions are retained at
reduced precision rather than evicted.  Holding the INT4 fraction at exactly
0.75 and the sink/recent protection common, role-aware placement and uniform
placement are indistinguishable: over 240 locked examples FullKV passes 213,
uniform-mixed 211, and structure-mixed 212, a one-example difference that no
paired test supports.  We report this as a negative result.

The implemented mixed BF16/INT4 path does reduce stored bytes, at a measured
cost.  Payload falls to 0.719$\times$ FullKV and allocated peak to
0.868$\times$, while time per output token rises by 1.20--1.215$\times$ and
end-to-end latency by 1.11--1.12$\times$.  The gap between the 0.719$\times$
measured payload and an idealised 0.473$\times$ byte model is accounted for by
storing each 4-bit code in a \texttt{uint8} container with FP32 scale and
zero-point metadata; the gap between payload and peak is accounted for by
dequantising cold pages into BF16 scratch for attention.  These are engineering
costs of the prototype rather than properties of the idea, but they are the
numbers this implementation achieves, and the latency comparison additionally
spans two different attention backends.

\begin{figure}[!htbp]
  \centering
  \includegraphics[width=0.86\linewidth]{systems_tradeoff.png}
  \caption{\textbf{Storage trade-off on H200.}  The implemented
  code-plus-metadata payload is smaller than the FullKV cache, while the
  allocated peak falls less because cold pages are expanded into BF16 scratch
  for attention.  Latency compares the mixed FlashInfer path against FullKV
  SDPA; the dashed line marks parity.}
  \label{fig:systems}
\end{figure}

\section{Discussion}
\label{sec:discussion}

\paragraph{A cheap signal with a knowable range.}
The practical value of a structural prior lies in what it does not require: no
attention statistics, no model behaviour, nothing beyond the prompt.  Our
results delimit where that suffices.  While the protected set fits the keep
budget, structural retention matches attention-based selection on the tasks we
can measure; past that point it degrades sharply and loses to it decisively.
Because $\rho$ is computable before generation, the choice between the free
signal and the paid one can be made per request.

\paragraph{But the ratio is not the whole diagnostic.}
The dissociation of Section~\ref{sec:sweep-mechanism} shows that $\rho$ predicts
whether the label can discriminate, not what is lost when it cannot.  A workload
whose critical state sits in the system preamble tolerates extreme
oversubscription; a workload whose critical state accumulates across turns does
not tolerate even modest oversubscription.  Agent traces are largely of the
second kind, which is an uncomfortable finding for the premise that motivated
this line of work, and it is the reason the same prior that succeeds on
schema-conformance tasks fails on BFCL.

\paragraph{On the value of soft weighting.}
ADAPT never loses to attention-based selection across the tested grid and
significantly exceeds it where that selection begins to degrade.  Because it
consumes attention scores, it cannot be justified as a cost reduction; its case
rests on accuracy at aggressive compression, and that case is currently
established on one model family in one regime.

\section{Limitations and Threats to Validity}
\label{sec:limitations}

\paragraph{The positive results are on a benchmark we authored.}
Every result showing an advantage is measured on PriorityBench-A.  The
externally authored evaluation shows the hard structural policy at 0.000, and
the pre-registered follow-up measured no ADAPT benefit there
($\Delta=-0.026$, $p=0.33$).  No claim of external superiority is made or
supported, and the reader should treat the ADAPT advantage of
Section~\ref{sec:sweep-adapt} as an in-distribution result whose external
counterpart has been sought and not found.

\paragraph{Instrument and stimulus are not independent.}
The tagger keys on role, length, and schema or constraint markers that resemble
the generator's visible templates, and the sweep's distractor schemas were
authored to carry roles the same tagger recognises.  The protected-mass axis is
measured post hoc with the pinned tokenizer and the mechanism replicates across
two model families, but this does not eliminate the concern, and it is the
principal reason the external evaluation rather than the sweep must carry any
method claim.

\paragraph{No difficulty gradient in the sweep.}
FullKV passes every cell of the sweep.  This isolates retention cleanly but
means the benchmark cannot speak to workload realism, and that pass rates are
close to binary.

\paragraph{Scope of the ADAPT result.}
The advantage over SnapKV is confined to Qwen3-8B at a 2\% budget.  On Llama the
comparison is uninformative rather than negative, because SnapKV does not
degrade enough to leave headroom.

\paragraph{External benchmark headroom.}
On BFCL, an 8B model passes only 20.5\% of conversations with the full cache and
fails 45.5\% with unparseable turn responses.  The workload therefore adjudicates
retention policies weakly, and our external conclusions are correspondingly
bounded.

\paragraph{Model, context, and budget coverage.}
Two 8B models, nominal contexts to 32k, and selected budgets between 2\% and
50\% are tested.  Behaviour at larger scales, longer contexts, and under batched
serving is unmeasured.  Systems measurements are single-request and do not
establish throughput, concurrency, or tail latency.

\paragraph{Relation to concurrent work.}
Role-aware retention is an active area, and we do not claim the use of
structural roles as a retention signal.  RoleKV~\cite{rolekv} predates this work
on the signal itself, and~\cite{structrolebias} predates it on the existence of
a density-dependent boundary, albeit at the sparse end of that axis and over
document-structural rather than conversational roles.  Our claims are confined
to the saturated regime, to lossy eviction, and to the tiebreak interaction of
Section~\ref{sec:sweep-mechanism}, which we believe to be new but which a reader
should weigh against that literature rather than in isolation.

\paragraph{Baseline fidelity.}
The sweep and external evaluations use the \texttt{kvpress} implementation of
SnapKV.  The earlier in-distribution comparisons additionally use repository
reimplementations of PyramidKV and a chunked H2O; the latter carries an explicit
fidelity caveat and should not be read as a reference-H2O result.

\section{Conclusion}
\label{sec:conclusion}

Application-visible structure is a retention signal that costs nothing to
compute, and this paper establishes the range over which it is sufficient.
Across 15{,}360 scored generations on two model families, structural retention
holds at ceiling while the protected set fits the keep budget and degrades
sharply beyond an oversubscription ratio of approximately one; past that point
it loses to matched-budget attention-based selection without winning a single
discordant pair.  The boundary is computable before generation, which makes it
usable as a serving-time diagnostic.

The more consequential finding is that the ratio alone is not that diagnostic.
Because the policy resolves ties by position, state carried in the leading system
block survives arbitrary oversubscription while state accumulated across
mid-conversation turns is lost entirely, and it is this second factor that
determines whether a workload is served by the cheap signal.  The distinction is
not a refinement of the earlier account but a correction of it, and it converts
an unexplained external failure into a predicted one.  That prediction was
registered before the corresponding experiment was scored and the experiment
returned it: on 233 paired conversations of an externally authored agent
benchmark, the blended policy is indistinguishable from attention-based
selection and trends slightly below it, as the mechanism requires.  The method
does not transfer.  The account of why it does not transfer does.

In operational terms: keep the schema before the pleasantry, but check first
whether the budget can hold the schema---and whether the state your task
actually depends on is the schema at all.

\begingroup
\small
\begin{thebibliography}{12}
\setlength{\itemsep}{0.45em}
\setlength{\parsep}{0pt}

\bibitem{vaswani2017attention}
Ashish Vaswani et al.
\newblock Attention is all you need.
\newblock \emph{Advances in Neural Information Processing Systems}, 30, 2017.

\bibitem{kwon2023pagedattention}
Woosuk Kwon et al.
\newblock Efficient memory management for large language model serving with PagedAttention.
\newblock In \emph{Proceedings of the 29th ACM Symposium on Operating Systems
Principles (SOSP)}, 2023.  \href{https://doi.org/10.1145/3600006.3613165}{doi:10.1145/3600006.3613165}.

\bibitem{xiao2024streamingllm}
Guangxuan Xiao et al.
\newblock Efficient streaming language models with attention sinks.
\newblock In \emph{International Conference on Learning Representations (ICLR)},
2024.  \href{https://openreview.net/forum?id=NG7sS51zVF}{OpenReview:NG7sS51zVF}.

\bibitem{zhang2023h2o}
Zhenyu Zhang et al.
\newblock H2O: Heavy-Hitter Oracle for efficient generative inference of large language models.
\newblock \emph{Advances in Neural Information Processing Systems}, 36, 2023.

\bibitem{li2024snapkv}
Yuhong Li et al.
\newblock SnapKV: LLM knows what you are looking for before generation.
\newblock \emph{Advances in Neural Information Processing Systems}, 37, 2024.
\href{https://doi.org/10.52202/079017-0722}{doi:10.52202/079017-0722}.

\bibitem{cai2024pyramidkv}
Zefan Cai et al.
\newblock PyramidKV: Dynamic KV cache compression based on pyramidal information funneling.
\newblock \href{https://arxiv.org/abs/2406.02069}{arXiv:2406.02069}, 2024.

\bibitem{liu2024kivi}
Zirui Liu et al.
\newblock KIVI: A tuning-free asymmetric 2bit quantization for KV cache.
\newblock In \emph{Proceedings of the 41st International Conference on Machine
Learning (ICML)}, PMLR 235:32332--32344, 2024.

\bibitem{ye2025flashinfer}
Zihao Ye et al.
\newblock FlashInfer: Efficient and customizable attention engine for LLM inference serving.
\newblock \emph{Proceedings of Machine Learning and Systems}, 7, 2025.

\bibitem{bai2024longbench}
Yushi Bai et al.
\newblock LongBench: A bilingual, multitask benchmark for long context understanding.
\newblock In \emph{Proceedings of the 62nd Annual Meeting of the Association for
Computational Linguistics}, pages 3119--3137, 2024.
\href{https://doi.org/10.18653/v1/2024.acl-long.172}{doi:10.18653/v1/2024.acl-long.172}.

\bibitem{hsieh2024ruler}
Cheng-Ping Hsieh et al.
\newblock RULER: What's the real context size of your long-context language models?
\newblock In \emph{First Conference on Language Modeling (COLM)}, 2024.
\href{https://openreview.net/forum?id=kIoBbc76Sy}{OpenReview:kIoBbc76Sy}.

\bibitem{patil2025bfcl}
Shishir G. Patil et al.
\newblock The Berkeley Function Calling Leaderboard (BFCL): From tool use to
agentic evaluation of large language models.
\newblock In \emph{Proceedings of the 42nd International Conference on Machine
Learning (ICML)}, PMLR 267:48371--48392, 2025.
\href{https://proceedings.mlr.press/v267/patil25a.html}{PMLR:patil25a}.

\bibitem{rolekv}
RoleKV: Role-aware KV cache management for the inverted age-importance of LLM
agent context.
\newblock OpenReview submission, 2026.
\href{https://openreview.net/forum?id=of1W47Odj1}{OpenReview:of1W47Odj1}.

\bibitem{structrolebias}
Soumil Mandal.
\newblock Adaptive filtering of the KV cache: Diagnosing and correcting
structural-role bias in LLM inference.
\newblock \href{https://arxiv.org/abs/2607.13205}{arXiv:2607.13205}, 2026.

\bibitem{valuediag}
Ruijie Zhang, Haozhe Liang, Da Chang, Li Hu, Fanqi Kong, Huaxiao Yin, and Yu Li.
\newblock When does value-aware KV eviction help?  A fixed-contract diagnostic
for non-monotone cache compression.
\newblock \href{https://arxiv.org/abs/2605.08234}{arXiv:2605.08234}, 2026.

\bibitem{yao2024taubench}
Shunyu Yao, Noah Shinn, Pedram Razavi, and Karthik Narasimhan.
\newblock $\tau$-bench: A benchmark for tool-agent-user interaction in real-world
domains.
\newblock In \emph{International Conference on Learning Representations (ICLR)},
2025.  \href{https://arxiv.org/abs/2406.12045}{arXiv:2406.12045}.

\end{thebibliography}
\endgroup
